\documentclass[11pt]{article}
\usepackage{amsmath,amssymb}
\title{A complete-cube signed incidence formula and a physical constant-flux test}\author{}\date{}
\begin{document}\maketitle
\paragraph{Input and scope.}
Use the signed gauge quotient and unsummed current calculation. The original $R=2$ cube of the full relative cube signed response calculation contains all 216
plaquettes, eight free links and seven independent chords. Its oriented
inverse-link insertion is BEFORE the inverse factor with a minus sign;
its positive-link insertion is AFTER the factor. The actual Cartan minimum
is odd in the boundary parameter, so its second derivative vanishes at
zero. Its exact seagull $77/6912$, bubble $1219/3456$, signed response
$-787/2304$, and full Schur energy $113/24$ agree with the incidence
formula below. The fixed-only curvature four remains in the classical
energy. This note extends the formula to every complete relative cube,
not the nonlinear estimate.

\paragraph{1. Original minimum and complete oriented incidence.}
Let $J$ be the non-tree free links of the original relative cube, with the
same rooted temporal forest as the relative green improved connected bound calculation. Fix a Cartan boundary path
$U_e(\epsilon)=\exp(i\epsilon b_e\sigma_3)$ and keep the tree links equal
to identity. Write the actual nearby minimum as
$U_j^*(\epsilon)=\exp(i\theta_j(\epsilon)\sigma_3)$.
Uniqueness near the flat minimum and invariance under global Cartan
conjugation imply this form. Conjugation by $i\sigma_1$ changes the sign
of the boundary parameter, so $\theta$ is odd. For each oriented plaquette
$p=((e_1,s_1),\ldots,(e_4,s_4))$ set
\[
 r_{pj}=\sum_{k:e_k=j}s_k,\quad K=r^*r,\quad G=K^{-1},\quad
 a_J=-G r^* d b,\quad f_p=\sum_{k=1}^4s_k a_{e_k}.                  \tag{1}
\]
Here $a=b$ on fixed boundary links and zero on the tree; $db$ means the
plaquette incidence of that boundary-only field. Then $a_J=\theta'(0)$,
$r^*f=0$, and $\theta''(0)=0$. The complete classical curvature is
$\sum_p f_p^2$, including fixed-only plaquettes. All these are actual
minimum derivatives; no fine field already integrated elsewhere is restored.

\paragraph{2. Neutral and charged Hessians with all inverse-link partners.}
Use right fluctuations at the actual minimum. The color-three Hessian
is a real matrix $N(\epsilon)$; the colors one and two form the real
representation of one Hermitian complex matrix $C(\epsilon)$. At zero
$N=C=K$. For two distinct chord occurrences $k<l$ in a plaquette define
\[
 \beta_{kl}=\sum_{k<q<l}s_q a_{e_q}
       +\mathbf1_{s_k=-1}s_k a_{e_k}
       +\mathbf1_{s_l=1}s_l a_{e_l},\qquad
 t_{kl}=f_p-2\beta_{kl}.                                            \tag{2}
\]
The endpoints in (2) are essential: positive variations occur after their
factor, negative variations before it. To second order in $\epsilon$, the
charged off-diagonal entry contributed by these two occurrences is
$s_ks_l\exp(-i\epsilon t_{kl})$, with the opposite entry its conjugate.
The immaterial overall imaginary sign changes with the quaternion convention.
The neutral contribution is $s_ks_l\cos(\epsilon f_p)$. Both diagonal
color blocks receive $\cos(\epsilon f_p)$ for every chord occurrence.
Consequently, with all repeated contributions summed,
\begin{align}
 N'(0)&=0,& N''(0)&=-r^*\operatorname{diag}(f_p^2)r,\\
 C'(0)&=iA,& A_{e_ke_l}&\mathrel{+}=-s_ks_l t_{kl},\quad
 A_{e_le_k}\mathrel{+}=s_ks_l t_{kl},\\
 C''_{jj}(0)&=-\sum_{p\ni j}f_p^2,&
 C''_{e_ke_l}(0)&\mathrel{+}=-s_ks_l t_{kl}^2.                       \tag{3}
\end{align}
In the last line the off-diagonal contribution is symmetric. $A$ is real
antisymmetric. These formulas follow by multiplying the ordered factors:
the transverse insertion pair has scalar part proportional to
$\cos(\epsilon(f_p-2\beta_{kl}))$ or its sine. The absence of
$\theta''(0)$ is why no missing minimum-acceleration term occurs at this
particular symmetric reference. For a general nonflat path it must be restored.

Since $\det H_{\rm real}=\det N\,(\det C)^2$, the FULL signed response is
\[
 \boxed{F_{\rm quad,\epsilon\epsilon}
 =\tfrac12\operatorname{Tr}(G N'')+\operatorname{Tr}(G C'')
                         +\operatorname{Tr}(G A G A).}             \tag{4}
\]
The last term is nonpositive; its negative is the positive covariance
bubble. The first two terms are the second-derivative partner, which need
not have a fixed sign. Equation(4) is the complete tree determinant, hence
already includes the gauge quotient of the signed gauge quotient and unsummed current calculation. Adding a separate ghost
factor to it would double count that quotient. The raw Gaussian normalizer
remains $e^{-V}$ times its original canonical/Haar factor times
$(\det H_{\rm real})^{-1/2}$; it is not the original nonlinear integral.

\paragraph{3. A physical direction with an exact Cartan minimum.}
On the complete cube of side $R$, choose on EVERY edge
\[
 a_{(x,\mu)}=\mathbf1_{\mu=1}\,{x_0\over R^2},                     \tag{5}
\]
and impose its boundary restriction as the boundary path. The temporal
forest (direction zero) is unchanged. Every 01 plaquette has
$f_p=R^{-2}$ and every other plaquette has zero flux. For every free link,
the two nonzero sine forces cancel exactly; all transverse color forces
vanish at the Cartan configuration. Thus (5) is the exact nearby action
minimum for sufficiently small $\epsilon$, not merely its linearized
harmonic extension. Its classical energy is
\[
 E_R=\sum_p f_p^2={(R+1)^2\over R^2},\qquad
 E_{R,\rm live}={(R-1)^2\over R^2},\qquad
 E_{R,\rm fixed}=4/R.                                               \tag{6}
\]
This is a genuine physical curvature direction, not a gauge orbit.
Its complete energy is bounded above and below as $R$ grows. The fixed-only
piece is part of the original exterior scalar; it cancels from the
classically-subtracted fluctuation but not from the complete physical energy.

\paragraph{4. Verified finite-size evidence, not an asymptotic estimate.}
The finite calculation builds ALL plaquettes and the exact original
forest, computes $G$, and implements (2)--(4). It reproduces
the full relative cube signed response calculation at $R=2$ to $10^{-12}$. For (5) it gives
\[
\begin{array}{c|r|r|r|r}
 R&|J|&\text{second-derivative term}&\text{positive bubble}&F''_{\rm quad}\\ \hline
 2&7&0.3489583333&0.4722222222&-0.1232638889\\
 3&80&6.1115422301&6.4315655746&-0.3200233445\\
 4&351&25.5232112541&26.0064691612&-0.4832579071
\end{array}
\]
The corresponding signed/complete-energy ratios are
$-0.0547839506,-0.1800131313,-0.3092850606$.
For canonical $\epsilon=\zeta/\sqrt B$, every reported fluctuation curvature
is divided by $B$, whereas the classical curvature is exactly $E_R$.
Large separate terms cancel substantially. Three sizes cannot determine
whether the combined response is logarithmic or establish any uniform
bound. Taking absolute values before combining (4) destroys precisely the
cancellation this test measures.

\paragraph{Remaining actual-model estimate.}
Equation(4) supplies a complete finite-incidence expression to which a
Green/curl estimate can now be applied without guessing insertion positions
or dropping gauge/acceleration partners. One must bound its COMBINATION
relative to (1)'s physical energy, including finite-boundary terms. Even
such a leading signed bound would not estimate the actual-law residual
(5) of the signed gauge quotient and unsummed current calculation. The latter still contains nonlinear current covariance and
actual second/fourth-moment differences. No uniform physical-scale,
coarse-iteration, or S94 conclusion is claimed.
\end{document}
